\documentclass[11pt]{article}

\usepackage[a4paper,margin=1in]{geometry}
\usepackage[T1]{fontenc}
\usepackage[utf8]{inputenc}
\usepackage{lmodern}
\usepackage{microtype}
\usepackage{setspace}
\usepackage{parskip}
\usepackage{hyperref}
\usepackage{enumitem}

\onehalfspacing

\title{Constitutional AI under Finite Capacity\\
Series Note}

\author{Panagiotis Kalomoirakis\\
Independent Researcher, Synkyria Project}
\date{April 2026}

\begin{document}

\maketitle

\begin{center}
\small\textsc{Canonical Release v1.0}
\end{center}

\begin{abstract}
This note explains the structure of the Constitutional AI under Finite Capacity series. The series is not organised as a single paper carrying every burden at once. Instead, it is organised as a linked set of six objects with distinct roles: a constitutional article draft, a note on constitutional force and applicability, an application profile, a minimal reviewable evidence specification, a clarification note on questions, limits, and misreadings, and a policy translation brief for public-governance audiences. Together, these objects define the constitutional core, its graded force, its projection into AI assistant systems, the minimum review surface required for constitutional claims, the disciplined limits of interpretation, and the public-facing translation of the proposal into governance language.
\end{abstract}

\section{Why this series is split}

A constitutional object for AI should not be forced to do every job at once. If a single paper tries to be constitution, explanatory bridge, application guide, evidence protocol, and critique response all together, it usually becomes either too abstract or too defensive.

The present series is therefore deliberately split into distinct objects. Each object carries a precise burden.

\section{The six objects}

\subsection*{1. The Finite-Capacity AI Constitution --- Article Draft}
This is the constitutional core. It states the minimal structural articles required for accountable AI operation under finite capacity.

\subsection*{2. Constitutional Force and Applicability in the Finite-Capacity AI Constitution}
This note explains how the Constitution should be read canonically: what functions as universal minimum constitutional core, what functions as stronger governance-level constitutional discipline, and what functions only where system class, domain profile, or runtime structure makes richer qualification necessary.

\subsection*{3. Minimal Application Profile for Constitutional AI in AI Assistant Systems}
This object projects the constitutional articles into a concrete system class: AI assistant systems.

\subsection*{4. Minimal Reviewable Evidence Specification for Constitutional AI}
This object states what evidence surfaces must exist if a constitutional claim is to be reviewable rather than merely asserted.

\subsection*{5. Questions, Limits, and Misreadings for Constitutional AI under Finite Capacity}
This clarification note prevents predictable misreadings without forcing the constitutional core itself to become a defensive text.

\subsection*{6. Constitutional AI under Finite Capacity: A Policy Translation Brief}
This brief translates the constitutional package into public-governance language for policy audiences, public authorities, compliance teams, standards-facing groups, and decision-makers.

\section{Recommended reading order}

The series should normally be read in the following order:
\begin{enumerate}[leftmargin=2em]
    \item \textbf{The Finite-Capacity AI Constitution}
    \item \textbf{Constitutional Force and Applicability in the Finite-Capacity AI Constitution}
    \item \textbf{Minimal Application Profile for Constitutional AI in AI Assistant Systems}
    \item \textbf{Minimal Reviewable Evidence Specification for Constitutional AI}
    \item \textbf{Questions, Limits, and Misreadings for Constitutional AI under Finite Capacity}
    \item \textbf{Constitutional AI under Finite Capacity: A Policy Translation Brief}
\end{enumerate}

This order matters.

The Constitution states the core articles.  
The Force and Applicability note explains how those articles are graded and how they should be read without overreach.  
The Application Profile shows how the constitutional grammar is projected into a concrete class of systems.  
The Evidence Specification states how constitutional claims become inspectable.  
The Questions note clarifies limits and prevents interpretive drift.
The Policy Translation Brief translates the package into public-governance language for external audiences.

For outward-facing use, the order of presentation may differ from the internal reading order. In public release settings, the Policy Translation Brief may appropriately be presented first as an entry document, followed by the Series Note or Instrument Index, and only then by the constitutional core and its companion instruments.

\section{How the objects relate}

The six objects should be read as a structured family rather than as isolated papers.

\begin{itemize}[leftmargin=2em]
    \item The \textbf{Constitution} states what must minimally hold.
    \item The \textbf{Force and Applicability note} states how strongly and under what conditions each part should be read.
    \item The \textbf{Application Profile} states where and how the constitutional grammar is projected.
    \item The \textbf{Evidence Specification} states how the constitutional claim becomes reviewable.
    \item The \textbf{Questions note} states what the series does not claim and how it should not be misread.
\end{itemize}

The \textbf{Policy Translation Brief} does not alter the constitutional core. Its function is different: it translates the package into public-governance language so that policy, regulatory, standards, and compliance audiences can understand the proposal without first entering the full internal architecture of the series.

The series therefore separates constitutional force, projection, review, and clarification into objects that can each remain precise.

\section{Historical continuity}

The series belongs to a broader programme in which constitutional structure, runtime evidence, and historically sensitive review are increasingly linked. Historical continuity does not function here as a universal demand for rich runtime machinery in every system. It functions more narrowly: where terminally similar verdicts are not semantically sufficient by themselves, relevant lineage or historical qualification must not be erased.

The series note does not itself define that principle in detail. It marks where in the series that burden is carried:
\begin{itemize}[leftmargin=2em]
    \item constitutionally, in the articles concerning qualification and lineage-sensitive interpretation;
    \item interpretively, in the note on constitutional force and applicability;
    \item and operationally, in richer runtime or evidence-bearing instantiations where historical non-equivalence matters.
\end{itemize}

\section{What this series is not}

This series is not:
\begin{itemize}[leftmargin=2em]
    \item a single implementation manual,
    \item a complete theory of ethics,
    \item a legal code,
    \item a finished compliance standard,
    \item or a claim that every AI system must instantiate the richest possible runtime architecture.
\end{itemize}

Its purpose is narrower and more disciplined: to provide a constitutional core, a graded reading of that core, a projection into a concrete system class, a minimal reviewable evidence layer, and a clarification surface for limits and misreadings.

\section{Closing statement}

The point of this series is not multiplication for its own sake. The point is structural clarity.

Constitutional AI under finite capacity becomes more legible when:
\begin{itemize}[leftmargin=2em]
    \item the constitutional core is kept distinct,
    \item its force and applicability are clarified separately,
    \item its application is projected through profiles,
    \item its claims are tied to reviewable evidence,
    \item and its limits are stated without overburdening the Constitution itself.
\end{itemize}

This series is therefore organised as a six-object architecture in which each document carries a precise burden and supports a clearer canonical reading of the whole.

\end{document}